\section{Pengujian Fungsional Aplikasi Web (Form, Alur, Respons)}

Pengujian fungsional memverifikasi bahwa setiap fitur bekerja sesuai spesifikasi. Untuk form, uji \textbf{skenario positif} (input valid menghasilkan output yang benar) dan \textbf{skenario negatif} (input tidak valid ditolak dengan pesan error yang jelas). Uji juga \textbf{boundary values} --- nilai di batas minimum dan maksimum. Untuk alur multi-langkah (misalnya registrasi $\to$ login $\to$ dashboard), pastikan setiap langkah berfungsi dan data mengalir dengan benar \cite{mdnToolsTesting}.

\begin{contoh}
\textbf{Skenario Pengujian Form Registrasi:}
\begin{itemize}
  \item \textbf{Positif}: isi semua field valid $\to$ data tersimpan, redirect ke login.
  \item \textbf{Negatif}: email tidak valid $\to$ pesan error ``Format email salah''.
  \item \textbf{Negatif}: password kurang dari 8 karakter $\to$ pesan error muncul.
  \item \textbf{Boundary}: nama tepat 1 karakter $\to$ apakah diterima atau ditolak?
  \item \textbf{Duplikasi}: email yang sudah terdaftar $\to$ pesan ``Email sudah digunakan''.
\end{itemize}
\end{contoh}

\begin{webcode}[language=PHP, caption={Validasi form registrasi untuk pengujian}]
<?php
$errors = [];
$nama  = trim($_POST['nama'] ?? '');
$email = trim($_POST['email'] ?? '');
$pass  = $_POST['password'] ?? '';

if (strlen($nama) < 3) $errors[] = "Nama minimal 3 karakter.";
if (!filter_var($email, FILTER_VALIDATE_EMAIL)) {
  $errors[] = "Format email tidak valid.";
}
if (strlen($pass) < 8) $errors[] = "Password minimal 8 karakter.";

if (empty($errors)) {
  // Simpan ke database dengan prepared statement
  echo "Registrasi berhasil.";
} else {
  foreach ($errors as $e) echo htmlspecialchars($e) . "\n";
}
?>
\end{webcode}

